\chapter{Literature Review}

\section{Overview of Healthcare Appointment Systems}

Healthcare appointment scheduling sits at the intersection of medicine, information technology, and service operations. Clinics must match finite doctor time with patient demand while minimising no-shows and idle capacity. Before computerisation, paper diaries and wall calendars performed this function. Computerised systems introduced searchable records, automated reminders, and statistical reports. Web-based systems further removed the need for patients to visit the clinic physically to reserve a slot. Cloud hosting removed the need for each clinic to maintain its own server room.

Digital health services have expanded rapidly. Patients expect to book consultations online, and providers need systems that respect privacy and role separation. For a computer science engineering project, it is essential to demonstrate not only user-facing features but also secure engineering practice: correct authentication, authorisation, and deployment on infrastructure that supports TLS encryption.

During my studies I identified a gap between simple demonstration websites and production-oriented applications. I therefore set out to build a complete system that could be demonstrated locally and also deployed on a public cloud URL with a real database, which is closer to industry practice than a prototype that runs only on a developer laptop.

\section{Theoretical Framework}

Healthcare appointment systems should be designed around the needs of real users rather than around database tables alone. User-Centered Design (UCD) provides a suitable theoretical basis for this thesis.

\subsection{User-Centered Design (UCD) Principles}

User-centered design is an iterative process that focuses on users and their needs at every stage of development. It emphasises understanding the target audience through research and design techniques so that the product is usable, accessible, and aligned with real tasks. The figure below shows the iterative principles of UCD (Figure~\ref{fig:ucd}).

\begin{figure}[htbp]
  \centering
  \includegraphics[width=0.95\textwidth]{figures/ucd-principles}
  \caption{UCD principles}
  \label{fig:ucd}
\end{figure}

In the context of \textbf{DocBook}, I did not run a full formal user-study programme with clinicians and patients at every iteration, but I applied UCD principles in a structured, intuitive way: I analysed how patients, doctors, and administrators work when appointments are booked manually, and I translated those needs into screens and workflows. For \textbf{patients}, the focus was on simple search, clear available slots, and visible appointment status. For \textbf{doctors}, the focus was on setting weekly availability, confirming bookings quickly, and recording prescriptions after a visit. For \textbf{administrators}, the focus was on overview dashboards and reports without exposing clinical tools meant for doctors.

This approach addressed typical problems in healthcare scheduling: long phone queues, unclear availability, double booking, and lack of a single view of the day's appointments. As noted in usability literature, understanding user needs and reflecting them in the interface is essential for adoption and correct use of digital health tools~\cite{iso9241}. DocBook was designed so that each role sees only relevant functions, which supports both usability and security.

\subsection{Core Features Guided by UCD Principles}

\begin{table}[htbp]
  \centering
  \caption{UCD principles mapped to DocBook features}
  \label{tab:ucd-features}
  \begin{tabular}{@{}p{4.5cm}p{8cm}@{}}
    \toprule
    UCD principle & DocBook feature \\
    \midrule
    Visibility of status & Appointment states: pending, confirmed, completed \\
    Error prevention & Unique constraint prevents double booking \\
    Recognition over recall & Dashboard lists upcoming appointments \\
    Role-appropriate UI & Separate patient, doctor, admin interfaces \\
    Trust and security & HTTPS, hashed passwords, RBAC \\
    \bottomrule
  \end{tabular}
\end{table}

\section{Comparative Analysis of Existing Systems}

\subsection{Healthcare and Appointment Platforms}

Commercial hospital systems integrate scheduling with billing and clinical records. They are powerful but expensive. Consumer booking platforms focus on discovery and booking for private practice. DocBook positions itself as an educational full-stack implementation with documented security and open deployability, not as a commercial competitor.

\begin{table}[htbp]
  \centering
  \caption{Comparison of approach types}
  \label{tab:approach-types}
  \begin{tabular}{@{}p{2.2cm}p{3.2cm}p{3.2cm}p{4.5cm}@{}}
    \toprule
    Type & Strength & Limitation & DocBook response \\
    \midrule
    Paper diary & Simple & No remote access & Web database \\
    Spreadsheet & Flexible & No RBAC & Multi-user app with roles \\
    Full HIS & Complete & High cost & Focused appointment module \\
    Generic demo & Fast build & No workflow & Full appointment lifecycle \\
    \bottomrule
  \end{tabular}
\end{table}

\subsection{Identified Gaps and Solutions}

\begin{table}[htbp]
  \centering
  \caption{Identified gaps and DocBook solutions}
  \label{tab:gaps}
  \begin{tabular}{@{}p{5.5cm}p{7cm}@{}}
    \toprule
    Gap & Solution \\
    \midrule
    No role separation & RBAC mixins per role \\
    Weak passwords & PBKDF2-SHA256 hashing \\
    No cloud HTTPS demo & Render deployment with TLS \\
    Double booking & Database uniqueness constraint \\
    No admin reporting & Appointment and revenue reports \\
    \bottomrule
  \end{tabular}
\end{table}
